\documentclass[11pt]{article}

\usepackage[margin=1in]{geometry}
\usepackage[T1]{fontenc}
\usepackage[utf8]{inputenc}
\usepackage{lmodern}
\usepackage{booktabs}
\usepackage{array}
\usepackage{hyperref}
\usepackage{xcolor}
\usepackage{enumitem}

\hypersetup{
    colorlinks=true,
    linkcolor=blue,
    urlcolor=blue
}

\title{Explainable IDS Pipeline: Execution Report and Documentation}
\author{}
\date{April 21, 2026}

\begin{document}
\maketitle

\section{Overview}
This document summarizes the purpose, workflow, fixes, and verified execution results of \texttt{pipeline.py}. The script implements an explainable intrusion detection system (IDS) workflow on the NSL-KDD dataset. It performs data loading, preprocessing, model training, evaluation, explanation generation, stability analysis, and adversarial analysis in a single end-to-end pipeline.

\section{Pipeline Scope}
The pipeline executes the following stages:
\begin{enumerate}[leftmargin=*]
    \item Load the NSL-KDD training and test sets.
    \item Convert the task to binary classification: \texttt{normal} vs.\ \texttt{attack}.
    \item One-hot encode categorical features and scale inputs with MinMax normalization.
    \item Train three classifiers:
    \begin{itemize}[leftmargin=1.5em]
        \item Logistic Regression
        \item Random Forest
        \item Multi-Layer Perceptron (MLP)
    \end{itemize}
    \item Evaluate each model using classification metrics and PR-AUC.
    \item Generate explanations with SHAP and LIME.
    \item Measure explanation stability across similar attack samples.
    \item Run a SHAP-guided adversarial evasion analysis and a simple defense evaluation.
\end{enumerate}

\section{Issues Found and Corrected}
The original runtime path had several execution problems. These were corrected directly in \texttt{pipeline.py}.

\subsection{Dataset Download Failure}
The original loader attempted to download the dataset exclusively from \texttt{raw.githubusercontent.com}. On this machine, DNS policy returned \texttt{NXDOMAIN} for that host, so the pipeline could not start from a clean workspace.

\textbf{Fix:}
\begin{itemize}[leftmargin=*]
    \item Kept the standard raw GitHub download path as the first attempt.
    \item Added an automatic fallback using the GitHub API blob endpoint.
    \item Added a clear, user-facing error if both download paths fail.
\end{itemize}

\subsection{Matplotlib Configuration Warning}
Matplotlib tried to write its configuration cache under a non-writable default path, producing warnings at startup.

\textbf{Fix:}
\begin{itemize}[leftmargin=*]
    \item Set \texttt{MPLCONFIGDIR} to a writable local directory under \texttt{outputs/.matplotlib}.
    \item Ensured the backend is configured before importing \texttt{pyplot}.
\end{itemize}

\subsection{Label Normalization Robustness}
The NSL-KDD family of files can appear with small formatting variations such as trailing periods in labels.

\textbf{Fix:}
\begin{itemize}[leftmargin=*]
    \item Normalized label strings with whitespace stripping.
    \item Treated both \texttt{normal} and \texttt{normal.} as the benign class.
\end{itemize}

\subsection{Explanation and Analysis Edge Cases}
Modern SHAP versions can return outputs in different binary-classification formats, and some downstream analysis steps can fail when no valid samples or pairs are available.

\textbf{Fix:}
\begin{itemize}[leftmargin=*]
    \item Added SHAP output normalization for list-based and ndarray-based return types.
    \item Added guards for empty attack subsets and empty similarity-pair results.
    \item Improved final summary handling when adversarial results are unavailable.
\end{itemize}

\section{Verified Execution}
The pipeline was run end-to-end successfully after the fixes. During execution, the missing dataset files were downloaded automatically through the GitHub API fallback.

\subsection{Dataset Sizes}
\begin{itemize}[leftmargin=*]
    \item \texttt{KDDTrain+.txt}: 19,109,424 bytes
    \item \texttt{KDDTest+.txt}: 3,441,513 bytes
\end{itemize}

\subsection{Observed Dataset Shapes}
\begin{itemize}[leftmargin=*]
    \item Training rows: 125,973
    \item Test rows: 22,544
    \item Final feature count after preprocessing: 122
\end{itemize}

\section{Model Results}
The final verified model metrics are shown in Table~\ref{tab:results}.

\begin{table}[h]
\centering
\begin{tabular}{>{\raggedright\arraybackslash}p{3.6in}cc}
\toprule
Model & Macro F1 & PR-AUC \\
\midrule
Logistic Regression (baseline) & 0.7535 & 0.8864 \\
Random Forest & 0.7726 & 0.9650 \\
MLP & 0.8060 & 0.9461 \\
\bottomrule
\end{tabular}
\caption{Verified evaluation results from the final pipeline run.}
\label{tab:results}
\end{table}

\noindent The MLP achieved the highest macro F1 score, while the Random Forest achieved the best PR-AUC.

\section{Explainability Results}
\subsection{SHAP}
The Random Forest SHAP analysis identified the following top features:
\begin{itemize}[leftmargin=*]
    \item \texttt{flag\_SF}
    \item \texttt{logged\_in}
    \item \texttt{dst\_host\_srv\_count}
    \item \texttt{dst\_host\_same\_srv\_rate}
    \item \texttt{same\_srv\_rate}
    \item \texttt{diff\_srv\_rate}
    \item \texttt{service\_private}
    \item \texttt{dst\_host\_rerror\_rate}
    \item \texttt{count}
    \item \texttt{dst\_host\_srv\_serror\_rate}
\end{itemize}

The Spearman rank correlation between Random Forest SHAP importances and MLP SHAP importances was:
\[
\rho = 0.745 \quad (p = 0.0000)
\]
This indicates substantial agreement between the two explanation profiles.

\subsection{LIME}
The top global LIME features from the evaluated attack instances were:
\begin{itemize}[leftmargin=*]
    \item \texttt{service\_private}
    \item \texttt{logged\_in}
    \item \texttt{wrong\_fragment}
    \item \texttt{dst\_host\_rerror\_rate}
    \item \texttt{rerror\_rate}
    \item \texttt{flag\_SF}
    \item \texttt{service\_ecr\_i}
    \item \texttt{dst\_host\_srv\_serror\_rate}
    \item \texttt{protocol\_type\_icmp}
    \item \texttt{flag\_S0}
\end{itemize}

\section{Stability and Adversarial Analysis}
\subsection{Explanation Stability}
For similar attack-sample pairs, the mean Jaccard similarity of the top-10 SHAP features was:
\[
0.994 \pm 0.037
\]
This indicates very high stability for the selected explanation subset under the implemented similarity criterion.

\subsection{SHAP-Guided Attack Results}
Table~\ref{tab:attack} summarizes the adversarial evasion rates.

\begin{table}[h]
\centering
\begin{tabular}{lccc}
\toprule
Perturbed Features & $\epsilon=0.05$ & $\epsilon=0.10$ & $\epsilon=0.20$ \\
\midrule
Top-5 & 34.86\% & 33.60\% & 32.26\% \\
Top-10 & 35.49\% & 35.04\% & 34.59\% \\
Top-15 & 35.49\% & 34.14\% & 33.51\% \\
\bottomrule
\end{tabular}
\caption{Verified evasion rates from the SHAP-guided perturbation analysis.}
\label{tab:attack}
\end{table}

\subsection{Defense Evaluation}
The implemented feature-randomization defense did not materially reduce evasion in the observed run:
\begin{itemize}[leftmargin=*]
    \item Evasion rate without defense: 35.04\%
    \item Evasion rate with defense: 35.04\%
    \item Observed reduction: 0.0\%
\end{itemize}

\section{Generated Artifacts}
The pipeline produced the following output figures in the \texttt{outputs/} directory:
\begin{itemize}[leftmargin=*]
    \item \texttt{shap\_beeswarm\_rf.png}
    \item \texttt{shap\_comparison.png}
    \item \texttt{shap\_stability.png}
    \item \texttt{evasion\_heatmap.png}
\end{itemize}

\section{How to Run}
From the project directory:

\begin{verbatim}
./venv/bin/python pipeline.py
\end{verbatim}

If the dataset files are not already present, the script will attempt to download:
\begin{itemize}[leftmargin=*]
    \item \texttt{KDDTrain+.txt}
    \item \texttt{KDDTest+.txt}
\end{itemize}

The downloader first tries the raw GitHub path, then falls back to the GitHub API if the raw content host is blocked.

\section{Conclusion}
The pipeline is now operational and verified on this machine. The main engineering change was making dataset acquisition resilient to host-specific DNS blocking of \texttt{raw.githubusercontent.com}. After that correction, the full IDS workflow executed successfully and generated model metrics, explainability outputs, stability measurements, and adversarial-analysis results without further runtime failures.

\end{document}
